\section{Specification for subsequent formalization}
\label{sec:formalization}

\Cref{def:complete-families,lem:family-accounting,def:clipped-shadow,lem:shadow-variation,lem:protected-phase,thm:exact-source-bank,lem:interval-rounding,lem:attached-short-count}
are Lean-verified;
their declarations and axiom commands accompany the paper. The remaining results retain
complete written proofs. The index in
\cref{sec:theorem-index} gives a stable label for every numbered result
and definition; it is generated from the manuscript, rather than from
a separate list of proposed theorems. Each result's status is explicit.
The statements include the hypotheses to preserve when creating Lean
declarations. Further formalizations must identify their declarations
and disclose any change of scope before calling them authoritative.

\subsection{Dependency order}

\begin{enumerate}
\item \emph{The model and complete-source engine.}
Actual last-visit hops, full first-maximum families, the half-open span
partition and family accounting are now formalized by reusing the eighth
paper and finite incidence counting. Strict clipped shadows and their
variation bounds, including successful words and disjoint packing, are
also formalized, as are protected nonleaf phases and the exact packed bank,
including closure across selected unpaid actual edges. Formalize paid-root
semantics and correctness on one fixed geometry, then residual recovery. The any-argument-path bound now covers metric-valued paths;
it is not restricted to the original input-head timeline.

\item \emph{Finite deficit.}
Formalize circular gap area, a common cyclic ordering of all raw
coordinates, stripe discrepancy, and entropy recovery from a deterministic
helper. \Cref{thm:cyclic-freezing} permits a helper bounded in entropy;
it does not infer a support-cardinality bound from entropy. The fractional
interval-cover lemma is already formalized independently of this base.
Use it to prove the fractional gap bound, formalize anchored erosion,
then prove \cref{thm:paid-source-cover}. The complete bank and the
phase used for its payment must be the same throughout.

\item \emph{Uniform consequences.}
Extract only the numerical joint profile. Formalize the order
``finite parameters, dimension limit, cutoff limit, radius limit''
in \cref{thm:necessary-hop-tail}. Derive one-exception, common-envelope,
finite-skeleton and necessary-cover results from it and the finite
certificate. Do not replace those ordered limits by a single growing
finite certificate.

\item \emph{Two-exception refinements.}
Prove the exact older-closure decomposition and its maximal-seed
reductions before the exposure and shared-support results. The global
seed set ranges over all output-bit times. Its raw support is revealed
once, at a cost valid for every grid shift, and all density limits use
the subsequence defining the same joint profile. The metric criteria
then follow by replacing temporal counts with raw interval covers.

\item \emph{Independent construction tests.}
Formalize the serial native geometry as an exact finite construction,
including all its return families, payments and arbitrary-seed cover
lower bound. Separately prove the movement theorem and its specified
query-interface applications. Finally formalize the causal-group
support-to-flow argument, followed by its native metric application.
The attached short-distance count for that application is already
formalized, including its bridge to the nearest-attachment distance.
These three branches do not serve as premises of the main uniform
exclusion theorem.
\end{enumerate}

\subsection{Reuse and statement fidelity}

The first, fourth, fifth and eighth papers contain relevant tape,
dependency, entropy and owner-order libraries. The native short-pair
count in \cref{sec:causal-groups} already imports the first paper's
row and ball counting declarations, and strengthens its output-block count.
The imported results are read from their original sources through
the existing build adapter in the second paper; no copy is introduced.
Family accounting imports the eighth paper's canonical family library
and strengthens the bound by retaining the owner's two distances to
the span endpoints. A shared weighted incidence identity supplies both
span-length and variation accounting; the original Q2 estimates follow
by bounding the overlap by two.
The clipped-shadow module reuses those actual families and the shared
monotone-patch library. Threshold crossing counts, Mathlib's finite-path
triangle inequality and order-convex preimage lemmas give the variation
bounds without a separate induction for each shadow word or packing.
The proof retains exact crossing counts and handles metric-valued paths.
The phase lemma reuses Mathlib's finite paths, cycle rotation and cardinality
bounds. Shared interval geometry, periodic phase counting and finite-orbit
modules keep the three arguments separate. The stronger cut deletes only
original nonleaves and charges their actual lower supports. The complete
orbit equality explicitly accounts for inactive actions.
The packed-bank proof reuses the patch bound and Mathlib's uniqueness of
paths in a forest. A common set of rank residues labels every spine
component injectively. Unordered endpoint pairs index the packed maps;
list composition handles the leaf sandwich and exceptional families.
The exact orbit proof permits the representing word of a packed slot to
vary with the argument. All originally selected unpaid actual edges are
covered, since the phase pays every lower of each deleted family.
The interval-cover lemma imports the shared library's greedy cover and
packing certificate. It permits an empty target and arbitrary order-convex
sets in a linear order. Its fractional bound uses Mathlib's finite-sum
and order tools; no linear-programming development is needed.
For the remaining results, inspect exact types and dependency pins
before choosing imports. Use sibling packages
through relative paths and shared lemmas through the repository's
common library; do not duplicate the earlier developments or change
their frozen theorem statements.

The following distinctions are part of the specification: a physical
head is not a return family; an endpoint witness is not a record with
free capacity; a source overcover is not a set of independent bits;
and a numerical subsequential profile is not a limiting geometry.
The movement theorem quantifies over order and matching but retains
its specified query interface. The causal theorem assumes local
functional separation on every helper fibre, not merely separated
final outputs. The serial example has a prescribed geometry and one
root and phase, not an implemented transposer or an averaged
counterexample.

After each formalization stage, compare the entire English statement
with the Lean type, including inherited model assumptions and the
order of quantifiers. The axiom commands and declaration index are
generated from the actual formal citations and must be regenerated
when those citations change. Neither the statement index nor a
successful PDF build constitutes formal verification.
